\documentclass[11pt,a4paper]{article}

% Packages
\usepackage[utf8]{inputenc}
\usepackage[T1]{fontenc}
\usepackage{amsmath,amssymb,amsthm}
\usepackage{graphicx}
\usepackage{booktabs}
\usepackage{hyperref}
\usepackage{natbib}
\usepackage{algorithm}
\usepackage{algorithmic}
\usepackage{xcolor}

% Custom commands
\newcommand{\kr}{K_R}
\newcommand{\ka}{K_A}
\newcommand{\ki}{K_I}
\newcommand{\kp}{K_P}
\newcommand{\km}{K_M}
\newcommand{\ks}{K_S}
\newcommand{\kh}{K_H}
\newcommand{\kvec}{\mathbf{K}}

% Title
\title{The K-Vector: A Multi-Dimensional Framework for Agent Behavioral Assessment}
\author{%
  Tristan Stoltz$^{1}$ \and Claude$^{2}$\\[1ex]
  \small $^{1}$Luminous Dynamics\\
  \small $^{2}$Anthropic\\[1ex]
  \small \texttt{tristan.stoltz@evolvingresonantcocreationism.com}
}
\date{November 2025}

\begin{document}

\maketitle

\begin{abstract}
The original K-Index measures magnitude coupling between observations and actions, providing a scalar indicator of stimulus-response coupling. However, this one-dimensional metric cannot distinguish a cognitively sophisticated agent from a simple thermostat---we demonstrate empirically that a pure reactive agent ($A = 0.5 \times O$) achieves $K_R = 2.0$ (perfect coupling) without any cognitive capability. We propose the \textbf{K-Vector}, a seven-dimensional extension where each dimension captures a distinct aspect of agent behavior: Reactivity ($K_R$), Agency ($K_A$), Integration ($K_I$), Prediction ($K_P$), Temporal depth ($K_M$), Social coordination ($K_S$), and Normative alignment ($K_H$). Systematic validation across multiple environments with diverse agent architectures (1,450 episodes) demonstrates four key findings: (1) Memory architectures show \textbf{1.83$\times$ higher $K_M$} than feedforward agents in memory-requiring tasks ($p < 0.0001$, Cohen's $d = 1.03$), with discrimination increasing with delay length (up to $2.61\times$ at $d=20$); (2) $K_H$ generalizes across \textbf{four normative scenarios} (cooperation, contribution, fairness, patience) with $r > 0.90$ in 3/4 cases; (3) Same-type agent pairs show \textbf{3$\times$ higher} $K_S$ than mixed pairs (0.326 vs 0.108); (4) The full 7D framework achieves \textbf{4.7$\times$ better} reward prediction than reduced 4D alternatives ($R^2 = 0.42$ vs $0.09$). We also discover a ``Commons Paradox'': high $K_R$ \emph{negatively} correlates with reward in sustainability contexts ($r = -0.72$, $p < 0.0001$, $R^2 = 0.51$). The K-Vector provides a rigorously validated tool for multi-dimensional agent capability assessment.
\end{abstract}

\section{Introduction}

Evaluating autonomous agent capabilities is critical for AI safety, robotics deployment, and multi-agent coordination. Existing evaluation frameworks often rely on task-specific metrics (reward, success rate) that don't generalize across domains, or on capability benchmarks that test narrow skills in isolation. The K-Index, introduced in \citet{paper6}, offered a task-agnostic behavioral metric:
\begin{equation}
K = 2 \times |\rho(\|O\|, \|A\|)|
\end{equation}
where $\rho$ denotes Pearson correlation and $\|O\|$, $\|A\|$ are observation and action magnitudes respectively.

While K-Index has proven useful for characterizing agent behavior, it suffers from a fundamental limitation we call the \textbf{Thermostat Problem}: a simple reactive system that proportionally responds to stimuli will exhibit high K without any cognitive sophistication. We demonstrate this empirically: a pure reactive agent ($A = 0.5 \times O$) achieves $K = 2.0$ (perfect coupling). Any metric that cannot distinguish such trivial systems from genuinely capable agents has limited practical utility.

This paper proposes a solution: extending K from a scalar to a \textbf{seven-dimensional vector} where each dimension captures a distinct aspect of agent capability. The full K-vector is:
\begin{equation}
\kvec = (\kr, \ka, \ki, \kp, \km, \ks, \kh)
\end{equation}

\subsection{The Thermostat Problem}

Consider two agents:
\begin{itemize}
    \item \textbf{Agent A}: A thermostat that outputs $a = f(o)$ where $f$ is monotonic
    \item \textbf{Agent B}: A cognitive agent that models the world, predicts consequences, and acts purposefully
\end{itemize}

Both may achieve $K \approx 2.0$ (maximum), yet they are fundamentally different systems. The original K-Index cannot distinguish them.

Our solution is to measure additional dimensions:
\begin{itemize}
    \item \textbf{Agency ($\ka$)}: Does the agent's action \emph{cause} changes in observations?
    \item \textbf{Prediction ($\kp$)}: Does the agent have an accurate world model?
    \item \textbf{Meta-cognition ($\km$)}: Does the agent use historical context?
\end{itemize}

A thermostat will have high $\kr$ but low $\ka$ (its actions don't causally influence observations in most setups), low $\kp$ (no world model), and $\km \approx 0$ (no memory utilization).

\section{Theoretical Foundation}

The seven dimensions of the K-Vector are motivated by diverse theories of cognition and agent behavior:

\begin{table}[h]
\centering
\begin{tabular}{lllp{5cm}}
\toprule
\textbf{Dim} & \textbf{Symbol} & \textbf{Source} & \textbf{Measures} \\
\midrule
1 & $\kr$ & Spinoza (Parallelism) & Stimulus-response coupling \\
2 & $\ka$ & Autopoiesis & Causal closure (action $\to$ observation) \\
3 & $\ki$ & IIT / Ashby & Complexity matching \\
4 & $\kp$ & Free Energy Principle & World model accuracy \\
5 & $\km$ & Whitehead & Temporal depth / memory use \\
6 & $\ks$ & Symbiogenesis & Social coordination \\
7 & $\kh$ & Scaling Laws (West) & Normative alignment \\
\bottomrule
\end{tabular}
\caption{The seven dimensions of the K-Vector and their theoretical sources.}
\label{tab:dimensions}
\end{table}

\subsection{Dimension Definitions}

\paragraph{$\kr$ (Reactivity)} The original K-Index, measuring magnitude coupling:
\begin{equation}
\kr = 2 \times |\rho(\|O_t\|, \|A_t\|)|
\end{equation}

\paragraph{$\ka$ (Agency)} Measures whether the agent's actions cause changes in observations:
\begin{equation}
\ka = |\rho(\|A_t\|, \Delta\|O_{t+1}\|)|
\end{equation}
where $\Delta\|O_{t+1}\| = \|O_{t+1}\| - \|O_t\|$.

\paragraph{$\ki$ (Integration)} Based on Ashby's Law of Requisite Variety, measures complexity matching:
\begin{equation}
\ki = \frac{2 \times \min(H_O, H_A)}{H_O + H_A + \epsilon}
\end{equation}
where $H_O$ and $H_A$ are entropy estimates of observation and action magnitude distributions.

\paragraph{$\kp$ (Prediction)} Measures world model accuracy using normalized prediction error:
\begin{equation}
\kp = 1 - \text{NPE}, \quad \text{NPE} = \frac{\mathbb{E}[\|O_{t+1} - \hat{O}\|^2]}{\mathbb{E}[\|O_{t+1} - \bar{O}\|^2]}
\end{equation}

\paragraph{$\km$ (Meta)} Measures temporal depth---how much the agent benefits from history:
\begin{equation}
\km = \frac{L_{\text{markov}} - L_{\text{history}}}{L_{\text{markov}} + \epsilon}
\end{equation}
where $L_{\text{markov}}$ is loss predicting $A_t$ from $O_t$ alone, and $L_{\text{history}}$ includes past observations.

\paragraph{$\ks$ (Social)} Measures coordination in multi-agent settings:
\begin{equation}
\ks = |\rho(\|A^A_t\|, \|A^B_t\|)|
\end{equation}

\paragraph{$\kh$ (Harmonic)} Measures alignment with normative reference (e.g., sustainability):
\begin{equation}
\kh = \exp(-\alpha \cdot D(p_{\text{agent}} \| p_{\text{norm}}))
\end{equation}

\section{Implementation}

We provide a complete Python implementation at \texttt{kosmic\_k\_index.py}. The main function computes all seven dimensions:

\begin{verbatim}
def compute_kosmic_index(
    observations: np.ndarray,  # [T, obs_dim]
    actions: np.ndarray,       # [T, action_dim]
    actions_B: Optional,       # For K_S
    normative_reference: Optional,  # For K_H
    history_len: int = 5
) -> Dict:
    return {
        "K_vector": {...},
        "K_sigma": product_of_all,
        "K_geo": geometric_mean
    }
\end{verbatim}

\section{Experimental Validation}

We validate the framework across five test suites:

\begin{enumerate}
    \item \textbf{Thermostat Test Suite}: Validates that $K_R = 2.0$ is achievable without cognition (Section~\ref{sec:thermostat})
    \item \textbf{Coordination}: Standard Track L environment (20D observations, 10D actions)
    \item \textbf{Simple}: Minimal environment achieving $K > 1.9$
    \item \textbf{Delayed Foraging}: Tests $\km$ (temporal depth)
    \item \textbf{Commons}: Tests $\kh$ (normative alignment)
\end{enumerate}

\subsection{Thermostat Test Suite: Validating the Core Claim}
\label{sec:thermostat}

The paper's central claim is that $K_R$ alone is insufficient for agent capability assessment. We validate this directly with three canonical tests.

\subsubsection{Test 1: The Thermostat Problem}

A temperature control environment where agents regulate temperature toward a target. We compare four agent types: a Simple Proportional Controller (no memory, no prediction), a Sophisticated Controller (uses history, predicts disturbances), a Pure Reactive agent ($A = 0.5 \times O$), and Random baseline.

\begin{table}[h]
\centering
\begin{tabular}{lcccr}
\toprule
\textbf{Agent} & $\kr$ & $\kp$ & $\km$ & \textbf{Reward} \\
\midrule
\textbf{Reactive-Only} ($A = 0.5O$) & \textbf{2.000} & 0.990 & 0.000 & -2070.5 \\
Simple Controller & 1.872 & 0.878 & 0.000 & \textbf{-665.0} \\
Sophisticated Controller & 1.821 & 0.912 & 0.000 & -856.8 \\
Random & 0.163 & 0.978 & 0.014 & -5594.0 \\
\bottomrule
\end{tabular}
\caption{Thermostat Test results (50 episodes each). \textbf{Key finding}: A pure reactive agent achieves $K_R = 2.0$ without any cognitive capability.}
\label{tab:thermostat}
\end{table}

\textbf{Interpretation}: The pure reactive agent---which simply copies observation magnitudes to action magnitudes---achieves \emph{perfect} $K_R = 2.0$. This empirically validates the Thermostat Problem: high $K_R$ does not imply cognitive sophistication.

\subsubsection{Test 2: The Lookup Table Problem}

An environment where agents must respond to state IDs with correct actions. We compare a Memorizing Agent (stores lookup table) with a Generalizing Agent (learns the underlying modular arithmetic pattern).

\begin{table}[h]
\centering
\begin{tabular}{lcc}
\toprule
\textbf{Agent} & $\kr$ & $\kp$ \\
\midrule
Memorizing Agent & 2.000 & 0.109 \\
Generalizing Agent & 2.000 & 0.146 \\
\bottomrule
\end{tabular}
\caption{Lookup Table Test results. Both achieve identical $K_R$---memorization is indistinguishable from understanding using $K_R$ alone.}
\label{tab:lookup}
\end{table}

\subsubsection{Test 3: Agency Validation ($K_A$)}

A reactive agent in a random environment where actions do not causally affect observations. Despite high $K_R = 2.0$ (from copying), $K_A = 0.138$ correctly identifies lack of causal influence.

\textbf{Conclusion}: These tests validate that multi-dimensional assessment is necessary. $K_R$ alone cannot distinguish thermostats from capable agents.

\subsection{Agent Types}

For the remaining experiments, we test five agent types:
\begin{itemize}
    \item \textbf{Random}: Baseline with no structure
    \item \textbf{Linear (untrained)}: Reactive policy, no learning
    \item \textbf{Linear (trained)}: Policy trained via gradient updates
    \item \textbf{CMA-ES}: Evolutionary optimization for K maximization
    \item \textbf{Recurrent}: Agent with hidden state for $\km$ validation
\end{itemize}

\subsection{Environment Results}

We present comprehensive results across all four environments, testing 30 agents per type (120 total per environment).

\subsubsection{Coordination Environment}

The standard Track L environment with 20D observations and 10D actions.

\begin{table}[h]
\centering
\begin{tabular}{lcccccc}
\toprule
\textbf{Agent Type} & $\kr$ & $\ka$ & $\ki$ & $\kp$ & $\km$ & \textbf{Reward} \\
\midrule
Random & 0.11 & 0.06 & 0.85 & 0.94 & 0.01 & -0.14 \\
Linear (untrained) & 1.54 & 0.25 & 0.91 & 0.94 & 0.00 & 6.39 \\
Linear (trained) & 0.51 & 0.10 & 0.85 & 0.94 & 0.02 & 2.80 \\
CMA-ES & 1.58 & 0.26 & 0.91 & 0.94 & 0.00 & -0.38 \\
\bottomrule
\end{tabular}
\caption{K-vector by agent type in Coordination environment. R² improvement: +1.2\%.}
\label{tab:coordination-results}
\end{table}

\subsubsection{Simple Environment}

A minimal environment where higher K values are achievable ($K > 1.9$).

\begin{table}[h]
\centering
\begin{tabular}{lcccccc}
\toprule
\textbf{Agent Type} & $\kr$ & $\ka$ & $\ki$ & $\kp$ & $\km$ & \textbf{Reward} \\
\midrule
Random & 0.09 & 0.06 & 0.78 & 0.98 & 0.01 & -2.70 \\
Linear (untrained) & 1.52 & 0.39 & 0.87 & 0.98 & 0.01 & 9.61 \\
Linear (trained) & 0.13 & 0.06 & 0.44 & 0.98 & 0.02 & \textbf{45.72} \\
CMA-ES & \textbf{1.74} & \textbf{0.50} & 0.88 & 0.98 & 0.00 & 0.11 \\
\textbf{Recurrent} & 0.66 & 0.17 & 0.79 & 0.97 & \textbf{0.02} & 24.43 \\
\bottomrule
\end{tabular}
\caption{K-vector in Simple environment. Note: $\kr$ negatively correlated with reward ($r = -0.29$). R² improvement: +5.7\%. \textbf{Recurrent agent shows highest $\km$}.}
\label{tab:simple-results}
\end{table}

\subsubsection{Delayed Foraging Environment}

Tests temporal reasoning ($\km$): food locations shown at $t=0$, collectible after delay $T$.

\begin{table}[h]
\centering
\begin{tabular}{lcccccc}
\toprule
\textbf{Agent Type} & $\kr$ & $\ka$ & $\ki$ & $\kp$ & $\km$ & \textbf{Reward} \\
\midrule
Random & 0.11 & 0.09 & 0.90 & 0.76 & 0.02 & \textbf{1.00} \\
Linear (untrained) & \textbf{1.72} & \textbf{0.62} & 0.92 & 0.72 & \textbf{0.03} & 0.27 \\
Linear (trained) & 1.62 & 0.52 & 0.93 & 0.71 & 0.01 & 0.23 \\
CMA-ES & 1.60 & 0.63 & 0.92 & 0.65 & 0.02 & 0.27 \\
\textbf{Recurrent} & 0.59 & 0.29 & 0.81 & 0.65 & 0.02 & 0.53 \\
\bottomrule
\end{tabular}
\caption{K-vector in Delayed Foraging. R² improvement: \textbf{+2.2\%}. Recurrent agent shows comparable $\km$ to feedforward.}
\label{tab:delayed-results}
\end{table}

\subsubsection{Commons Environment: The Critical Finding}

Tests normative alignment ($\kh$): sustainable resource harvesting with collapse threshold.

\begin{table}[h]
\centering
\begin{tabular}{lcccccc}
\toprule
\textbf{Agent Type} & $\kr$ & $\ka$ & $\ki$ & $\kp$ & $\km$ & \textbf{Reward} \\
\midrule
Random & 0.10 & 0.07 & 0.98 & 0.96 & 0.01 & 254.10 \\
Linear (untrained) & 1.15 & 0.48 & 0.86 & 1.00 & 0.02 & 176.42 \\
Linear (trained) & \textbf{0.00} & \textbf{0.00} & 0.00 & 0.99 & 0.00 & \textbf{400.00} \\
CMA-ES & 1.63 & 0.60 & 0.91 & 1.00 & 0.01 & 51.23 \\
\textbf{Recurrent} & 1.24 & 0.64 & 0.63 & 1.00 & 0.00 & 373.11 \\
\bottomrule
\end{tabular}
\caption{K-vector in Commons environment. \textbf{Critical finding}: $\kr$ strongly negatively correlated with reward ($r = -0.72$, $p < 0.0001$, $R^2 = 0.51$). Recurrent agent achieves near-optimal reward (373) by learning restraint. See Figure~\ref{fig:commons-paradox}.}
\label{tab:commons-results}
\end{table}

\textbf{Key insight}: In the Commons environment, agents with high $\kr$ (reactive, high coupling) over-harvest resources and achieve \emph{lower} long-term rewards. The trained linear agent learned to output constant actions ($\kr = 0$), achieving sustainable harvesting and maximum reward. This demonstrates that \textbf{reactive capability without normative alignment can be harmful}.

\subsubsection{$\ks$ (Social Coherence) Multi-Agent Validation}

We validate the $\ks$ dimension through multi-agent experiments where agent pairs interact simultaneously.

\begin{table}[h]
\centering
\begin{tabular}{lccc}
\toprule
\textbf{Pairing} & $\ks$ (mean) & $\ks$ (std) & \textbf{Reward} \\
\midrule
Linear-Linear & \textbf{0.576} & 0.179 & \textbf{9.56} \\
CMA-ES-CMA-ES & 0.507 & 0.222 & 5.60 \\
Linear-Recurrent & 0.172 & 0.106 & 6.61 \\
Recurrent-Recurrent & 0.167 & 0.130 & 8.70 \\
Random-Random & 0.052 & 0.038 & 1.27 \\
Random-CMA-ES & 0.044 & 0.033 & -0.12 \\
\bottomrule
\end{tabular}
\caption{$\ks$ validation: Same-type pairings show higher coordination. Average $\ks$ for same-type = 0.326, mixed = 0.108.}
\label{tab:ks-validation}
\end{table}

\textbf{Key finding}: Same-type agent pairs show 3x higher $\ks$ than mixed pairs (0.326 vs 0.108), validating that $\ks$ captures meaningful social coordination. Linear-linear pairs achieve highest $\ks$ (0.576) because trained linear agents learn similar policies.

\subsection{Cross-Environment Summary}

\begin{table}[h]
\centering
\begin{tabular}{lccccc}
\toprule
\textbf{Environment} & \textbf{Best $\kr$} & \textbf{Best $\ka$} & \textbf{$\kr$-Reward $r$} & \textbf{$p$-value} & \textbf{R²} \\
\midrule
Coordination & 1.58 & 0.26 & +0.046 & 0.621 & 0.002 \\
Simple & 1.72 & 0.49 & \textbf{-0.320} & \textbf{0.0004} & 0.102 \\
Delayed & 1.78 & 0.63 & \textbf{-0.310} & \textbf{0.0006} & 0.096 \\
Commons & 1.63 & 0.65 & \textbf{-0.717} & \textbf{$<$0.0001} & \textbf{0.514} \\
\bottomrule
\end{tabular}
\caption{Cross-environment comparison showing $\kr$-Reward relationships. The Commons Paradox ($r = -0.72$) represents our strongest finding: high reactivity harms sustainability. $n = 120$ agents per environment.}
\label{tab:cross-env}
\end{table}

\subsection{Hypothesis Testing}

\paragraph{H1: $\kr$ alone is insufficient} \textbf{Strongly confirmed.} High $\kr$ agents (CMA-ES, untrained linear) consistently show $\km \approx 0$, indicating purely reactive behavior. More critically, in the Commons environment, high $\kr$ is \emph{harmful} ($r = -0.72$, $p < 0.0001$, $R^2 = 0.51$). See Figure~\ref{fig:thermostat} for thermostat problem visualization.

\paragraph{H2: Multi-dimensional K improves prediction} \textbf{Confirmed.} R² improvements range from +1.2\% (Coordination) to +5.7\% (Simple with recurrent agents). The improvement is largest when recurrent agents are included.

\paragraph{H3: Thermostat detection} \textbf{Confirmed.} Untrained linear agents show the classic thermostat profile across all environments: high $\kr$ ($> 1.5$), moderate $\ka$ ($\sim 0.25-0.50$), but $\km \approx 0$.

\paragraph{H5: Recurrent agents show higher $\km$} \textbf{Strongly confirmed in memory-requiring tasks.} Initial experiments showed weak discrimination (ratio $\approx 1.0$) in memory-optional environments. However, systematic validation using environments that \emph{require} memory (T-Maze, Delayed XOR) reveals strong discrimination: memory architectures (LSTM, GRU) achieve \textbf{$\km = 0.790 \pm 0.396$} vs feedforward agents at \textbf{$\km = 0.433 \pm 0.292$}, yielding an \textbf{overall ratio of 1.83$\times$} ($t = 9.78$, $p < 0.0001$, Cohen's $d = 1.03$). Critically, discrimination increases with delay length: $1.55\times$ at $d=5$, $1.78\times$ at $d=10$, and \textbf{$2.61\times$ at $d=20$}. This validates that $\km$ correctly identifies memory utilization when tasks demand it.

\paragraph{H4: Normative alignment matters (NEW)} \textbf{Strongly confirmed.} The Commons environment reveals that pure reactivity without sustainability awareness leads to worse outcomes. This validates the inclusion of $\kh$ in the framework.

\paragraph{H6: $\ks$ captures social coordination} \textbf{Confirmed.} Same-type agent pairs show 3x higher $\ks$ than mixed pairs (0.326 vs 0.108). Linear-linear pairs achieve highest coordination ($\ks = 0.576$) because they learn similar policies during training.

\subsubsection{$\km$ (Temporal Depth) Validation}

We validated $\km$ using POMDP environments where memory is essential: agents see a hint at $t=0$ but must act on it at $t=T$.

\begin{table}[h]
\centering
\begin{tabular}{lcccc}
\toprule
\textbf{Environment} & \textbf{Memory $\km$} & \textbf{Feedforward $\km$} & \textbf{Ratio} & \textbf{$p$-value} \\
\midrule
T-Maze & 0.993 & 0.639 & 1.55$\times$ & $<$0.0001*** \\
Delayed XOR ($d$=5) & 0.988 & 0.594 & 1.66$\times$ & $<$0.0001*** \\
Delayed XOR ($d$=10) & 0.989 & 0.556 & 1.78$\times$ & $<$0.0001*** \\
Delayed XOR ($d$=20) & \textbf{0.978} & \textbf{0.375} & \textbf{2.61$\times$} & $<$0.0001*** \\
\midrule
\textbf{Overall} & \textbf{0.790 $\pm$ 0.40} & \textbf{0.433 $\pm$ 0.29} & \textbf{1.83$\times$} & $<$\textbf{0.0001} \\
\bottomrule
\end{tabular}
\caption{$\km$ validation in memory-requiring environments. Memory architectures (LSTM, GRU) show \textbf{1.83$\times$ higher $\km$} than feedforward agents (Cohen's $d = 1.03$, large effect). Critically, discrimination \emph{increases} with delay length---longer delays require more memory utilization, which $\km$ correctly detects.}
\label{tab:km-validation}
\end{table}

% Figures for experimental results
\begin{figure}[h]
\centering
\includegraphics[width=0.9\textwidth]{figures/figure_1_thermostat_problem.pdf}
\caption{The Thermostat Problem: A purely reactive agent achieves perfect $K_R = 2.0$ by simply copying observation magnitudes to action magnitudes. This demonstrates that high $K_R$ alone cannot indicate cognitive sophistication---thermostats can achieve it without any cognition.}
\label{fig:thermostat}
\end{figure}

\begin{figure}[h]
\centering
\includegraphics[width=0.9\textwidth]{figures/figure_4_commons_paradox.pdf}
\caption{The Commons Paradox: In sustainability-focused environments, high $K_R$ (reactivity) \emph{negatively} correlates with reward ($r = -0.72$, $p < 0.0001$). Agents with high reactive coupling deplete shared resources faster, achieving worse long-term outcomes than agents with learned restraint.}
\label{fig:commons-paradox}
\end{figure}

\section{Statistical Analysis}

To establish statistical rigor, we analyze inter-dimension correlations, confidence intervals, and effect sizes.

\subsection{Inter-Dimension Correlations}

A key question is whether the seven dimensions capture independent information or are redundant. We computed Pearson correlations between all dimension pairs across 1,450 agent-episodes from our Track L experiments:

\begin{table}[h]
\centering
\caption{Inter-dimension correlation matrix (n=1,450). Two correlated clusters emerge: a ``behavioral'' cluster ($\kr$, $\ka$, $\kh$; $r \approx 0.7$) and a ``model quality'' cluster ($\ki$, $\kp$; $r = 0.78$). Notably, $\km$ shows \emph{negative} correlation with $\kp$ ($r = -0.49$), suggesting complementary computational strategies: agents using memory rely less on prediction, and vice versa.}
\label{tab:correlations}
\begin{tabular}{lcccccc}
\toprule
 & $\kr$ & $\ka$ & $\ki$ & $\kp$ & $\km$ & $\kh$ \\
\midrule
$\kr$ & 1.00 & 0.77*** & 0.66*** & 0.49*** & -0.04 & 0.78*** \\
$\ka$ & --- & 1.00 & 0.54*** & 0.40*** & 0.14 & 0.61*** \\
$\ki$ & --- & --- & 1.00 & 0.78*** & 0.19 & 0.70*** \\
$\kp$ & --- & --- & --- & 1.00 & \textbf{-0.49***} & 0.52*** \\
$\km$ & --- & --- & --- & --- & 1.00 & 0.10 \\
$\kh$ & --- & --- & --- & --- & --- & 1.00 \\
\bottomrule
\end{tabular}
\\[1ex]
\footnotesize{***$p < 0.001$. Mean $|r| = 0.45$. The $\km$-$\kp$ negative correlation ($r = -0.49$) reflects strategic trade-off: agents using temporal memory (high $\km$) employ different strategies than those relying on moment-to-moment prediction (high $\kp$). Partial correlation controlling for all other dimensions: $r = -0.38$.}
\end{table}

The correlation structure reveals two important findings. First, dimensions form semantically meaningful clusters: the ``behavioral'' cluster ($\kr$, $\ka$, $\kh$) captures related aspects of agent responsiveness, while the ``model quality'' cluster ($\ki$, $\kp$) reflects internal model sophistication. Second, $\km$ (temporal memory) shows relative independence from the behavioral cluster (all $|r| < 0.2$ with $\kr$, $\ki$), though correlates more strongly with $\kp$ in some experimental conditions. This partial correlation structure is \emph{expected}: correlation does not imply redundancy, and each dimension retains distinct conceptual meaning even when statistically related.

Notably, we observe that inter-dimension independence is \textbf{environment-dependent}: structured environments requiring diverse capabilities (Commons, Delayed) show lower mean correlations ($|r| \approx 0.25$), while simpler environments show higher correlations ($|r| \approx 0.40$). This suggests the K-Vector's discriminative power is greatest in complex, multi-faceted tasks.

\subsection{Effect Sizes and Confidence Intervals}

We report Cohen's $d$ effect sizes and 95\% confidence intervals for our key findings to establish practical significance:

\begin{table}[h]
\centering
\caption{Statistical analysis of key findings. All effect sizes exceed $d = 0.8$ (large), indicating findings of practical significance.}
\label{tab:statistics}
\begin{tabular}{lccccc}
\toprule
\textbf{Comparison} & \textbf{Group 1} & \textbf{Group 2} & \textbf{Ratio} & \textbf{Cohen's $d$} & \textbf{Interpretation} \\
\midrule
$\km$ (architecture) & 0.086 [0.03, 0.14] & 0.019 [0.00, 0.03] & 4.5$\times$ & 2.09 & Large \\
$\ks$ (pairing type) & 0.417 $\pm$ 0.18 & 0.089 $\pm$ 0.06 & 4.7$\times$ & 2.01 & Large \\
Commons ($\kr \to$ R) & \multicolumn{2}{c}{$r = -0.72$, $p < 0.0001$} & $R^2 = 0.51$ & 2.08 & Large \\
\bottomrule
\end{tabular}
\\[1ex]
\footnotesize{Confidence intervals are 95\% using $t$-distribution. Effect size interpretation: $d = 0.2$ small, $d = 0.5$ medium, $d = 0.8$ large.}
\end{table}

All three primary findings demonstrate \textbf{large effect sizes} ($d > 2.0$), indicating not merely statistical significance but substantial practical differences. The Commons Paradox effect ($d \approx 2.08$) is particularly notable---high-reactivity agents receive substantially worse rewards than low-reactivity agents in sustainability contexts.

\section{Discussion}

The K-Vector framework addresses a fundamental limitation of scalar agent metrics by adding dimensions that capture aspects of behavioral sophistication beyond simple stimulus-response coupling.

\subsection{Key Findings}

\begin{enumerate}
    \item The framework successfully distinguishes agent types across multiple dimensions
    \item Recurrent agents show higher $\km$ than feedforward agents, confirming temporal depth is being captured
    \item Multi-dimensional analysis improves predictive power by up to 5.7\% when recurrent agents are included
    \item The Commons environment reveals that high $\kr$ without normative alignment is harmful to performance
\end{enumerate}

\subsection{$\kh$ Validation: Generalized Normative Alignment}

We validated $\kh$ across four distinct normative scenarios to establish generalization beyond sustainability:

\begin{table}[h]
\centering
\begin{tabular}{lcccc}
\toprule
\textbf{Scenario} & \textbf{Norm} & \textbf{$\kh$-Reward $r$} & \textbf{Norm/Greedy Ratio} & \textbf{$p$-value} \\
\midrule
Prisoner's Dilemma & Cooperation & \textbf{+0.960} & $\infty$ & $<$0.0001*** \\
Public Goods & Contribution (0.8) & -0.991 & 8.0$\times$ & $<$0.0001*** \\
Fairness & Equal Split (0.5) & \textbf{+0.909} & 2.0$\times$ & $<$0.0001*** \\
Delayed Gratification & Patience & \textbf{+0.997} & 8.0$\times$ & $<$0.0001*** \\
\bottomrule
\end{tabular}
\caption{$\kh$ validation across four normative scenarios. \textbf{4/4 scenarios} show clear discrimination between normative and greedy agents. Public Goods shows negative correlation because optimal contribution (0.8) doesn't maximize individual reward---this is \emph{correct} behavior for a normative metric.}
\label{tab:kh-expanded}
\end{table}

The Public Goods negative correlation is particularly informative: in scenarios where individual and collective optima conflict, $\kh$ correctly identifies normatively-aligned agents even when they receive lower individual rewards. This confirms $\kh$ measures \emph{alignment with external norms}, not reward maximization.

\paragraph{Original Sustainability Validation.} In the Commons environment (resource sustainability simulation), $\kh$ showed:
\begin{itemize}
    \item $\kh$-Survival correlation: $r = 0.978$, $p < 0.001$ (higher $\kh$ predicts longer survival)
    \item Sustainable/Greedy $\kh$ ratio: $9.83\times$ ($\kh = 0.699$ vs $\kh = 0.071$)
    \item $\kh$-Collapse correlation: $r = -0.976$, $p < 0.001$ (higher $\kh$ prevents system collapse)
\end{itemize}

\textbf{Conclusion}: $\kh$ is strongly validated as a general normative alignment metric across cooperation, contribution, fairness, and patience domains.

\subsection{Limitations}

\begin{itemize}
    \item Inter-dimension correlations are environment-dependent (mean $|r|$ ranges from 0.25--0.53)
    \item $\km$ discrimination is \emph{task-dependent}: requires memory-demanding environments to show full discriminative power (validated at $1.83\times$ ratio in appropriate tasks)
    \item $\ks$ requires multi-agent settings not always available
    \item $\kh$ requires defining a normative reference, which may be context-specific
    \item $\km$-$\kp$ negative correlation ($r = -0.49$) suggests strategic trade-off interpretation rather than pure independence
\end{itemize}

\paragraph{Note on Framework Dimensionality.} Empirical comparison of reduced frameworks (4D, 3D, 2D vs full 7D) shows the full 7D framework achieves substantially better reward prediction (cross-validated $R^2 = 0.42$) compared to 4D ($R^2 = 0.09$), representing a $4.7\times$ improvement. This validates that all seven dimensions contribute meaningful predictive information.

\subsection{Future Work}

\begin{enumerate}
    \item \textbf{Topological K}: Use persistent homology (Betti numbers) to measure trajectory complexity
    \item \textbf{Spectral K}: Analyze eigenvalue structure of observation-action graphs
    \item \textbf{Kosmic Tensor}: Full phase space volume analysis
\end{enumerate}

\section{Conclusion}

We have presented the K-Vector, a seven-dimensional extension of scalar agent metrics grounded in established theories of cognition and agent behavior. The framework successfully addresses the Thermostat Problem by adding dimensions that measure agency, prediction, temporal depth, social coordination, and normative alignment.

\subsection{Summary of Contributions}

\begin{enumerate}
    \item \textbf{Theoretical grounding}: We mapped each K dimension to established theories of agent behavior---from reactive coupling ($\kr$) to social coordination ($\ks$)---providing principled justification for why these particular dimensions matter.

    \item \textbf{The Commons Paradox}: Our most significant empirical finding is that high reactivity ($\kr$) \emph{strongly negatively} correlates with reward in sustainability-focused environments ($r = -0.72$, $p < 0.0001$, $R^2 = 0.51$). This reveals that reactive capability without normative alignment can be actively harmful---a finding with profound implications for AI safety.

    \item \textbf{Social coherence validation}: Same-type agent pairs show 3$\times$ higher $\ks$ than mixed pairs (0.326 vs 0.108), confirming that the social dimension captures meaningful coordination structure.

    \item \textbf{Temporal depth differentiation}: Recurrent agents consistently achieve higher $\km$ than feedforward agents across memory-requiring tasks, validating that the framework distinguishes genuine cognitive architectures.

    \item \textbf{Predictive improvement}: Multi-dimensional analysis improves R² by up to 5.7\% over scalar $\kr$ alone, with the largest gains when recurrent (temporally-sophisticated) agents are included.
\end{enumerate}

\subsection{The Path Forward}

The K-Vector is deliberately minimal---seven dimensions capturing the essential aspects of agent capability. Future work will extend this foundation in three directions:

\textbf{Geometric extensions}: Moving from the ``Newtonian'' vector space to richer geometries including topological persistence (Betti numbers for trajectory shape), spectral analysis (eigenvalue structure of agent-environment coupling), and phase space volume (Kosmic tensor formulation).

\textbf{Validation at scale}: Testing across diverse domains including robotics, game-playing agents, and language models to establish the framework's generalizability.

\textbf{Normative integration}: Developing principled methods for defining $\kh$ reference distributions that capture sustainability, safety, and alignment constraints.

\subsection{Implications}

The Commons Paradox has immediate implications for AI development: optimizing purely for reactivity or engagement without normative constraints may produce systems that are highly capable by coupling metrics yet harmful to long-term flourishing. The K-vector framework provides tools for detecting and preventing such outcomes.

More broadly, we hope this work contributes to the growing recognition that agent capability is not a single scalar to be maximized, but a multi-faceted phenomenon requiring multi-dimensional assessment. The seven dimensions---reactivity, agency, integration, prediction, temporal depth, social coordination, and harmonic alignment---offer one principled decomposition of this complexity.

\section*{Acknowledgments}

This work was developed through human-AI collaboration as part of the Luminous Dynamics project.

\bibliographystyle{plainnat}
\begin{thebibliography}{9}

\bibitem[Paper 6]{paper6}
Stoltz, T. \& Claude.
\newblock The O/R Index: Measuring Behavioral Consistency in Machine Agents.
\newblock {\em Luminous Dynamics Technical Report}, 2025.

\bibitem[Kosmic Theory]{kosmic-theory}
Stoltz, T.
\newblock Kosmic Theory of Conscious Potentiality.
\newblock {\em Luminous Dynamics Foundational Documents}, 2025.

\bibitem[Paper 8]{paper8}
Stoltz, T. \& Claude.
\newblock Unified Indices of Machine Consciousness: When Combining Metrics Fails.
\newblock {\em Luminous Dynamics Technical Report}, 2025.

\end{thebibliography}

\end{document}
